\documentclass[12pt,a4paper]{article}  % Déclare la classe du document.

%  ------------------------------------------  Modules usuels  --------------------------------------
%  --------------------------------------------------------------------------------------------------------

\usepackage{amsmath} 	
\usepackage{amssymb} 
\usepackage{listings}         
\usepackage{numprint} 
\usepackage{pifont}

%  --------------------------------------------------------------------------------------------------------
% ----------------------  packages pour fontes et caractères spéciaux  ---------------------
%  --------------------------------------------------------------------------------------------------------

\usepackage{eurosym}   
\usepackage{amsfonts}   
\usepackage{amstext}       
\usepackage{amsfonts}	
\usepackage{mathrsfs}	 

\usepackage{pstricks,pst-all,pstricks-add,pst-tree,pst-node,pst-3dplot}

%------------------------------------------------------------------------------------------------------
%  --------------------------  Francisation et gestion des caractères sous mac --------------------------------------------------------------------------------------------------------

\usepackage[francais]{babel}   
\usepackage[utf8]{inputenc}
\usepackage[OT1,T1]{fontenc} 
%  --------------------------------------------------------------------------------------------------------
% ------------------------------------Packages  de  mise en page ---------------------------
%--------------------------------------------------------------------------------------------------------

\usepackage{fancyhdr}	% en-têtes et pieds de pages
\usepackage{amsfonts}

 \usepackage{multicol} % permet l'utilisation du multicolonnage dans le texte. La syntaxe est \begin{multicols}{nb de colonnes}    ......  \end{multicols}
 
\usepackage{calc}            % permet de faire un minimum de calcul sur les arguments des commandes et macros (utilisé dans la def de l'environnement preuve )
\usepackage{epic}            % extension et simplification de  l'environnement picture
\usepackage{graphicx}    % permet d'inclure des graphique externe dans un document
\usepackage{variations}   % permet d'obtenir de joli tableau de variations
\usepackage{enumerate}  %permet de faire des listes personnalisées
\usepackage{arydshln}     %permet d'avoir des lignes en pointillés
\usepackage{framed}  
\usepackage{multido} 


%  --------------------------------------  MACROS USUELLES  ----------------------------------------

\newcommand{\R}{\mathbb{R}} 
\newcommand{\N}{\mathbb{N}} 
\newcommand{\D}{\mathbb{D}} 
\newcommand{\Z}{\mathbb{Z}} 
\newcommand{\Q}{\mathbb{Q}} 
\newcommand{\C}{\mathbb{C}} 

%  -------------------------------------  MACROS SPECIFIQUES ------------------------------------


\newcommand{\raisedrule}[2][0em]{\leaders\hbox{\rule[#1]{1pt}{#2}}\hfill}


\newcounter{numexo}

\setcounter{numexo}{1}

\newcommand{\exercice}[1]{\vskip 0.6cm\noindent\rule[+2pt]{1cm}{0.5pt}\ {\large\bf \ Exercice \thenumexo {  #1 }}\raisedrule[+2pt]{0.5pt}\null\addtocounter{numexo}{1}\par\vskip 0,4cm}


\newcommand{\lignes}[1]{
\setlength{\unitlength}{1mm}
\vskip 0,4cm\noindent
\multido{\i=0+1}{#1}{%
     \noindent
     \begin{picture}(150,5)
     \linethickness{0,005mm}
     \put(-5,0){\dashline{1}(175,5)(5,5)}
     \end{picture}\vskip 0,001cm
     }
\vskip 0cm
}




% ----------------------------------------- FORMAT DE PAGE  -----------------------------------------

\usepackage[left = 1.3cm,
	 	      right = 1.3cm,
		      top = 1.7cm,
		      headsep = 8mm,
		      bottom = 1.7cm,
		      headheight = 14.2 pt
		      ]{geometry}

\setlength{\parindent}{0mm}


\everymath{\displaystyle} 
\renewcommand{\thesection}{\Roman{section}} 
\author{}
\title{}
\date{}
\pagestyle{fancy}

 
\lhead{1ère - spécialité maths} 


\rhead{\textcolor{red}{Correction - DS n$^{\circ}$7} - 18/05/21}		


\fancyfoot[C]{} 
\fancyfoot[L]{} 
\fancyfoot[R]{} 

\lstset{language=Python}  

\lstset{frame = single
framerule = 1pt
showspaces = false,
showstringspaces = false,
numbers=left, numberstyle=\tiny}

% ----------------------------------------- correction  -----------------------------------------
\usepackage[most,skins,breakable]{tcolorbox}
\usepackage{enumitem}
\usepackage{tikz,tkz-tab}

\newcommand{\corr}[1]{\fcolorbox{red}{yellow!30}{#1}}

\newcommand{\bareme}[1]{\tikz\node[circle,draw,inner sep=1pt,outer sep=1pt, color=red]{\textbf{#1}};}

\newtcolorbox{corrBox}{enhanced,boxrule=1pt,nobeforeafter,boxsep=0pt,left=\fboxsep,right=\fboxsep,colback=yellow!30,colframe=red!75!black}

\newlist{enumMet}{enumerate}{2}
\setlist[enumMet,1]{label=\arabic*), font=\bf\color{blue!75!black}}
\setlist[enumMet,2]{label=\alph*), font=\bf\color{blue!75!black}}


\begin{document}
\graphicspath{{images/}}

\begin{center}

\fbox{\LARGE{\rule{0.5cm}{0pt}\rule[-0.25cm]{0pt}{0.9cm} \textcolor{red}{Correction - DS N$^{\circ}$7 - Application de la dérivation	}\rule{0.5cm}{0pt}}}
\end{center}


% \begin{center}\textbf{1h - Calculatrice autorisée - Barème indicatif}\end{center}

% \begin{center}\textbf{Les élèves avec un tier-temps ne traitent pas les questions avec  le symbole \ding{46}}\end{center}

% \begin{center}
% %On veillera à ENCADRER ses résultats et soigner sa copie.\\
% Les résultats doivent être justifiés par des calculs (au moins 1 étape intermédiaire).
% \end{center}
\vspace{.05cm} 

\vspace{-0.3cm}
\exercice{ - Tableau de variations - ( 10 points )}

Pour chaque fonction, définie sur I, Calculer sa dérivée, puis dresser son tableau de variations
\begin{enumMet}
    \item $f: x \mapsto 3x^{2}-6 x+9, \mathrm{I}=\mathbb{R}$ \quad \quad \bareme{2,5}\\
    \begin{corrBox}
        $f'(x)=6x-6$. $f'$ est une fonction affine avec $a=6>0$. La fonction est donc croissante et négative puis positive. Le signe change quand $x=1$\\
        \begin{tikzpicture}
        \tkzTabInit{$x$ / 0.8 , $f'(x)$ / 0.8,$f(x)$ / 1.5}{$-\infty$, $1$, $+\infty$}
        \tkzTabLine{, -,z ,+,  }
        \tkzTabVar{+/  ,-/$6$ , +/ }
    \end{tikzpicture}\\
    \end{corrBox}
    \item $g: x \mapsto x^{3}-6 x^{2}+9 x-1, \mathrm{I}=\mathbb{R}$\quad \quad \bareme{2,5}\\
    \begin{corrBox}
        $g'(x)=3x^2-12x+9$. \\
        $\Delta=36 \iff x_1=1 \text{ et }x_2=3$. De plus, $a=3>0$\\
        \begin{tikzpicture}
        \tkzTabInit{$x$ / 0.8 , $f'(x)$ / 0.8,$f(x)$ / 1.5}{$-\infty$, $1$,$3$, $+\infty$}
        \tkzTabLine{, +,z ,-,z,+,  }
        \tkzTabVar{-/  ,+/$3$ , -/$-1$, +/ }
    \end{tikzpicture}\\
    \end{corrBox}
    \item $h: x \mapsto \dfrac{3 x-1}{x+2}, I=\mathbb{R} \backslash\{-2\}$.\quad \quad \bareme{2,5}\\
    \begin{corrBox}
        $h'(x)=\dfrac{3(x+2)-1(3x-1)}{(x+2)^2}=\dfrac{7}{(x+2)^2}>0$. \\
        \begin{tikzpicture}
        \tkzTabInit{$x$ / 0.8 , $f'(x)$ / 0.8,$f(x)$ / 1.5}{$-\infty$, $-2$, $+\infty$}
        \tkzTabLine{, +,d ,+,  }
        \tkzTabVar{-/  ,+D-/ , +/ }
    \end{tikzpicture}\\
    \end{corrBox}
    \item \ding{46} $k: x \mapsto x^{4}-\frac{4}{3} x^{3}, \mathrm{I}=\mathbb{R}$\quad \quad \bareme{2,5}\\
    \begin{corrBox}
        $k'(x)=4x^3-4x^2=4x^2(x-1)$. \\
        $k'$ est du signe de $(x-1)$ car $4x^2>0$ pour $x\in\R$\\
        \begin{tikzpicture}
        \tkzTabInit{$x$ / 0.8 , $f'(x)$ / 0.8,$f(x)$ / 1.5}{$-\infty$, $0$,$1$, $+\infty$}
        \tkzTabLine{, -,z ,-,z,+,  }
        \tkzTabVar{+/  ,R/ , -/$-\dfrac{1}{3}$, +/ }
        \tkzTabIma{1}{3}{2}{$0$}
    \end{tikzpicture}\\
    \end{corrBox}
\end{enumMet}

\exercice{ - extremums locaux - ( 5 points )}
On suppose que $f$ et $g$ sont des fonctions définies et dérivables sur un intervalle donné. Déterminer les éventuels extremums locaux des fonctions $f$ et $g$ sur cet intervalle.
\begin{enumMet}
    \item $f(x)=2x(3x^2-18)$\quad \quad \bareme{2,5}\\
    \begin{corrBox}
        $f'(x)=2(3x^2-18)+6x(2x)=18x^2-36=18(x^2-2)=18(x-\sqrt{2})(x+\sqrt{2})$\\
        \begin{tikzpicture}
        \tkzTabInit{$x$ / 1 , $f'(x)$ / 1,$f(x)$ / 2}{$-\infty$, $-\sqrt{2}$,$\sqrt{2}$, $+\infty$}
        \tkzTabLine{, +,z ,-,z,+,  }
        \tkzTabVar{-/  ,+/$24\sqrt{2}$ , -/$-24\sqrt{2}$, +/ }
        \end{tikzpicture}\\
         Le minimum local sur $\R$ est $-30$ et le maximum local sur $\R$ est $30$
    \end{corrBox}
    \item $g(x)=4x-5+\dfrac{3}{x-2}$\quad \quad \bareme{2,5}\\
    \begin{corrBox}
        $g'(x)=4-\dfrac{3}{(x-2)^2}=\dfrac{4(x-2)^2-3}{(x-2)^2}=\dfrac{4x^2-16x+13}{(x-2)^2}$\\
        Le signe de g dépend de celui du numérateur. On étudie le signe du numérateur. \\
        $4x^2-16x-13=0\iff \Delta=48\iff x_1=\dfrac{-\sqrt{3}+4}{2} \text{ et } x_2=\dfrac{\sqrt{3}+4}{2}$\\
        \begin{tikzpicture}
        \tkzTabInit{$x$ / 1 , $f'(x)$ / 1,$f(x)$ / 2}{$-\infty$, $x_1$,$2$,$x_2$, $+\infty$}
        \tkzTabLine{, +,z ,-,d,-,z,+,  }
        \tkzTabVar{-/  ,+/$-4\sqrt{3}+3$ ,-D+/, -/$4\sqrt{3}+3$, +/ }
        \end{tikzpicture}\\
        Pas de minimum et pas de maximum.
    \end{corrBox}
\end{enumMet}

\exercice{ - Recette et bénéfice - ( 6 points )}

Une entreprise fabrique et vend des trottinettes électriques. On modélise le coût total de fabrication journalier, en centaines d'euros, de $x$ trottinettes ($2\leqslant x\leqslant 14)$ par la fonction notée $C$ telle que :
$$
C(x)=\frac{1}{3} x^{3}-6 x^{2}+36 x-50
$$
Le nombre $x$ de trottinettes est fonction de leur prix unitaire $P$, qui vérifie  la condition $P=12-\frac{1}{2} x$.

\begin{enumMet}
    \item \begin{enumMet}
        \item Déterminer la recette journalière $R(x)$ correspondant à la vente de $x$ trottinettes.\quad \quad \bareme{1}\\
        \begin{corrBox}
            $R(x)=x\times P(x)=x\left(12-\dfrac{1}{2}x\right)=12x-\dfrac{1}{2}x^2$
        \end{corrBox}
        \item Dresser le tableau de variations de la fonction $R(x)$\quad \quad \bareme{1}\\
        \begin{corrBox}
            $R'(x)=-x+12$. $R'$ est une fonction affine avec $a=-1<0$. La fonction est donc décroissante et positive puis négative. Le signe change quand $x=12$\\
            \begin{tikzpicture}
                \tkzTabInit{$x$ / 0.8 , $R'(x)$ / 0.8,$R(x)$ / 1.5}{$2$, $12$, $14$}
                \tkzTabLine{, +,z ,-,  }
                \tkzTabVar{-/$22$  ,+/$72$ , -/$70$ }
            \end{tikzpicture}\\
        \end{corrBox}
        \item Pour quel nombre de trottinettes vendues quotidiennement réalise-t-on la recette maximale ? Quel est dans ce cas le prix d'une trottinette?\quad \quad \bareme{1}\\
        \begin{corrBox}
            On réalise la recette maximale(7200\EUR{}) pour 12 trottinettes vendues quotidiennement. Dans ce cas le prix d'une trottinette est $P=\left(12-\dfrac{1}{2}\times12\right)\times100=600$\EUR{}
        \end{corrBox}
    \end{enumMet}
    \item \begin{enumMet}
        \item Exprimer le bénéfice journalier $B(x)$ obtenu pour la vente de $x$ trottinettes.\quad \quad \bareme{1}\\
        \emph{Rappel : Le bénéfice d'une entreprise se calcule par différence entre la recette et le coût de production.}\\
        \begin{corrBox}
            $B(x)=R(x)-C(x)=12x-\dfrac{1}{2}x^2-\left(\dfrac{1}{3} x^{3}-6 x^{2}+36 x-50\right)=-\dfrac{1}{3} x^{3}+\dfrac{11}{2}x^2-24x+50$
        \end{corrBox}
        \item \ding{46} Dresser le tableau de variations de la fonction $B(x)$\quad \quad \bareme{1}\\
        \begin{corrBox}
            $B'(x)=-x^2+11x-24$\\
            $B'(x)=0\iff \Delta=25\iff x_1=3 \text{ et } x_2=8$\\
            \begin{tikzpicture}
                \tkzTabInit{$x$ / 0.8 , $B'(x)$ / 0.8,$B(x)$ / 1.5}{$2$, $3$,$8$, $14$}
                \tkzTabLine{, -,z ,+,z,-,  }
                \tkzTabVar{+/$21{,}3$  ,-/$18{,}5$,+/$39{,}3$ , -/$-122$ }
            \end{tikzpicture}\\
        \end{corrBox}
        \item \ding{46} Pour quel nombre de trottinettes vendues réalise-t-on le bénéfice journalier maximal ? Quel est dans ce cas le prix d'une trottinette électrique?\quad \quad \bareme{1}\\
        \begin{corrBox}
            Il faut vendre 8 trottinettes pour réaliser le bénéfice maximal ($3900$\EUR{}). Dans ce cas le prix d'une trottinette est $P=\left(12-\dfrac{1}{2}\times8\right)\times100=800$\EUR{}
        \end{corrBox}
    \end{enumMet}
    % \item \textbf{Question Bonus}
    % \begin{enumMet}
    %     \item Avec la calculatrice, représenter dans un même repère les courbes représentatives des fonctions $C$, $R$ et $B$ sur [2; 14].
    %     \item Que peut-on dire du bénéfice journalier de cette entre prise lorsque la recette journalière est maximale? Justifier.
    % \end{enumMet}
\end{enumMet}

\end{document}

